% !TeX program = pdflatex
% !TeX encoding = UTF-8
% P0：面向碎石堆小行星任务的锚定机构多指标评价与证据化比较框架。
% 当前状态：研究规划；尚未完成系统检索、数据提取或机构评分。
% 本篇的贡献不能只是主观加权表格，必须包含任务条件、硬约束、证据质量、
% 不确定性传播和排名稳定性。若文献证据不足，形成“不可稳定排序”的定量结论
% 与标准化测试协议同样有效。
%
% 候选命题：
% H1 不同任务场景会实质改变锚定机构排序。
% H2 只比较最大承载力会误判全过程任务效用。
% H3 将证据不确定性传播到评分后，部分表面优势会变得不可分辨。
% H0 一组简单工程指标已足以稳定排序；复杂框架没有新增决策信息。
%
% 图路线：F01 技术分类；F02 成功事件链；F03 指标/硬约束；F04 证据图谱；
% F05 场景评分；F06 不确定性和权重；F07 性能—证据；F08 研究优先级。

\documentclass[11pt]{article}
\usepackage[T1]{fontenc}
\usepackage[utf8]{inputenc}
\usepackage{amsmath,amssymb}
\usepackage{booktabs}
\usepackage{graphicx}
\usepackage[hidelinks]{hyperref}
\usepackage[margin=1in]{geometry}

\title{A Mission-Conditioned and Evidence-Aware Framework\\
for Evaluating Asteroid Anchoring Systems}
\author{Authors to be confirmed}
\date{Planning manuscript --- evidence pending}

\newcommand{\planned}[1]{\par\noindent\textit{Planned work: #1}\par}
\newcommand{\figureplaceholder}[1]{%
  \fbox{\parbox{0.91\linewidth}{\centering
  \vspace{0.018\textheight}\textbf{Planned evidence --- no results yet}\par
  \smallskip #1\par\vspace{0.018\textheight}}}}
\newif\ifBibliographyReady
\BibliographyReadyfalse

\begin{document}
\maketitle

\begin{abstract}
This planning manuscript develops a mission-conditioned framework for comparing
asteroid anchoring systems using process-level success, resource demand, terrain
compatibility, and evidence uncertainty. The proposed method separates hard
feasibility constraints from compensatory utilities and propagates incomplete or
uncertain literature evidence into scenario-specific scores and rankings. No
systematic review, mechanism score, or technology ranking has yet been completed.
\end{abstract}

\section{Decision problem and scope}
% 科学问题：为何现有承载力数据不能直接支持任务选型？评价对象是“具体配置+条件”，
% 不是给整类技术贴固定分数。界定小行星锚定，不无限扩展到所有地外基础工程。
% 完成条件：至少冻结三类任务、评价对象、系统边界和决策使用者。
\planned{Define representative mission decisions, the anchoring process boundary,
and the configurations to compare. Distinguish short contact, long-duration
surface operation, and high-load drilling or excavation scenarios.}

\begin{figure}[htbp]
\centering
\figureplaceholder{P0-F01: Taxonomy of impact penetrators, drilled anchors,
microspines, force-closure systems, expandable subsurface anchors, adhesive or
solidification concepts, capillary concepts, and staged hybrids.}
\caption{Planned technology taxonomy. Each evidence record will refer to a
specific configuration and test condition rather than an entire mechanism class.}
\end{figure}

\section{Mission success model}
% 区分：地表适用/接触、部署、载荷转移、任务期保持；各概率通常相关。
% 不能把“试验成功率”和长期可靠度重复加分。失败定义、任务时间和重试规则要统一。
Let the mission-level success event be decomposed conditionally as
\begin{equation}
P_{\mathrm{mission}}=
P(C)P(D\mid C)P(T\mid D,C)P(H\mid T,D,C),
\end{equation}
where $C$, $D$, $T$, and $H$ denote compatible contact, deployment, load
transfer, and load retention, respectively. The conditional representation does
not imply independent stages.

\planned{Define measurable events, failure states, retry rules, mission duration,
and how correlated terrain uncertainty propagates through the event chain.}

\begin{figure}[htbp]
\centering
\figureplaceholder{P0-F02: Event tree from terrain encounter to sustained task
load, including abort, retry, release, and unrecoverable failure.}
\caption{Planned mission-success event chain and conditional probabilities.}
\end{figure}

\section{Metrics, hard gates, and utility}
% 硬约束：任务力/矩包络、安装反力、环境、时间、质量/体积/功率/耗材上限。
% 连续性能：成功率、承载裕度、资源、扰动、适应性、复用性等。
% 避免相关指标重复计权；效用函数体现达标后的边际收益，而非盲目线性归一化。
For mission scenario $j$, define a feasibility gate
\begin{equation}
G_j=\prod_{r=1}^{R}\mathbb{I}\!\left[g_{r,j}(\boldsymbol{x})\leq 0\right],
\end{equation}
and a conditional utility
\begin{equation}
U_j=\sum_{i=1}^{I}w_{i,j}u_{i,j}(x_i),
\qquad \sum_i w_{i,j}=1.
\end{equation}
A nominal score is $S_j=100G_jU_j$. Final reporting will propagate input,
model, and missing-data uncertainty rather than presenting this point value alone.

\planned{Elicit defensible utility curves and weights from mission requirements;
separate non-compensatory gates from graded preferences; test double-counting and
the consequences of alternative aggregation rules.}

\begin{figure}[htbp]
\centering
\figureplaceholder{P0-F03: Metric hierarchy, hard gates, conditional utilities,
and links to traceable mission requirements.}
\caption{Planned evaluation structure. A high score in one metric cannot rescue
a configuration that fails a hard mission constraint.}
\end{figure}

\section{Systematic evidence collection}
% 工作：检索方案、时间范围、数据库、纳入排除、双人/复核提取、重复出版物处理。
% 记录载荷方向、自由度、地表材料、重力、真空、台架、样本量、失败及测量误差。
% 固定台架抗拔不能直接代表自由飞行器全过程锚定。
\planned{Pre-register search, screening, extraction, and harmonization rules.
Record geometry, terrain, gravity, boundary support, load direction, sample size,
failures, and uncertainty for every usable observation. Preserve missingness.}

\begin{figure}[htbp]
\centering
\figureplaceholder{P0-F04: Evidence map showing which mechanisms have concept,
simulation, terrestrial, vacuum, reduced-gravity, or flight evidence.}
\caption{Planned evidence coverage and missing-data map.}
\end{figure}

\section{Uncertainty and evidence-aware comparison}
% 不把成熟度随意乘到性能上。分别报告潜在性能、证据强度、保守决策分数。
% 方案：概率输入/区间 + 蒙特卡洛；报告中位数、5%分位、区间、越门概率。
% 权重与效用曲线做全局敏感性，识别稳定排序、场景专用和不可区分方案。
\planned{Propagate aleatory and epistemic uncertainty to gate satisfaction and
utility. Report score distributions, conservative quantiles, rank probabilities,
and sensitivity to weights, utility shapes, priors, and missing-data assumptions.}

\section{Planned comparative results}
% 结果未完成。不能预写某类方案最佳。对照：不同任务场景、不同证据处理、
% 只看承载力的简单排名、完整全过程框架。分析帕累托前沿与支配关系。
\begin{figure}[htbp]
\centering
\figureplaceholder{P0-F05: Scenario-specific score distributions and feasibility
probabilities for configurations supported by comparable evidence.}
\caption{Planned mission-conditioned comparison; no ranking is established yet.}
\end{figure}

\begin{figure}[htbp]
\centering
\figureplaceholder{P0-F06: Rank stability under weights, uncertainty, missing data,
and alternative aggregation rules.}
\caption{Planned robustness and sensitivity analysis.}
\end{figure}

\begin{figure}[htbp]
\centering
\figureplaceholder{P0-F07: Performance potential versus evidence strength,
with uncertainty and untested domains shown explicitly.}
\caption{Planned separation of potential performance from evidential maturity.}
\end{figure}

\begin{figure}[htbp]
\centering
\figureplaceholder{P0-F08: Value-of-information or priority map identifying
measurements that most reduce decision uncertainty.}
\caption{Planned research and standardization priorities.}
\end{figure}

\section{Validation, interpretation, and limits}
% 方法验证：已知支配关系、极限情况、合成数据；外部专家/需求一致性；
% 不用所比较方案的结果反过来任意调权重。若排名不稳，结论是不可区分。
\planned{Test limiting cases, internal consistency, and synthetic recovery.
Separate scientific evidence from stakeholder preference and report when data do
not support a stable ranking. Propose common test protocols only where gaps are
demonstrated by the evidence map and sensitivity analysis.}

\section{Completion and stop conditions}
% 成篇：系统检索+可审计数据+任务指标+硬门槛+不确定性+权重敏感性+可复用工具。
% 止损：若数据异质性使任何量化映射不可辩护，则转为系统证据图谱与标准协议；
% 若只能做主观表格，则作为项目内部工具，不单独投稿。
\planned{Proceed as an independent paper only if the framework yields auditable,
uncertainty-aware comparisons or a rigorous demonstration of non-comparability.
A subjective weighted checklist alone is insufficient.}

\section*{Data and software availability}
\planned{Archive the screened-record table, data dictionary, scoring code,
scenario definitions, and versioned sensitivity settings before submission.}

\ifBibliographyReady
  \bibliographystyle{plain}
  \bibliography{references}
\else
  \section*{References to verify}
  Verified sources will be added after the search protocol and comparison scope
  are frozen. The bibliography is intentionally disabled at the planning stage.
\fi

\end{document}
